% !TeX root = ../../Book3.tex
\section{价公式与模形式环}
\label{b3:sec:2305}

普通辐角原理在基本区域上仍有效，但圆弧配对时函数出现权因子，椭圆固定点又只贡献部分圆周。价公式就是把这两种修正逐项记录的结果。

\subsection{基本域中的零点}
\label{b3:subsec:230501}

对非恒零亚纯模形式，记 \(v_p(f)\) 为上半平面局部参数中的零点阶数，极点取负数。模群变换中的因子在上半平面非零，故阶数只依赖轨道。取每条普通轨道一个代表；\(\mathrm{i}\) 与 \(\rho\) 单列。尖点阶数按 \(q\) 参数定义。

\ProseParagraphBreak

固定点的阶数还有同余限制。例如 \(f(-1/\tau)=\tau^kf(\tau)\) 在 \(\mathrm{i}\) 附近比较首项，得到 \((-1)^{v_i}=\mathrm{i}^k\)。权为二时必在 \(\mathrm{i}\) 消失，足以排除非零权二全纯形式；一般同余关系也可以精化维数计数。

\subsection{价公式}
\label{b3:subsec:230502}

\begin{theorem}[价公式]\label{b3:thm:23-valence}
若 \(f\) 是非零权 \(k\) 亚纯模形式，在尖点也亚纯，则
\[
v_\infty(f)+\frac12v_{\mathrm{i}}(f)+\frac13v_\rho(f)
+\sum_{p\ne\mathrm{i},\rho}v_p(f)=\frac{k}{12},
\]
其中求和按普通模群轨道计数。
\end{theorem}
\begin{proof}
先假定边界没有零极点。在高度 \(Y\) 截断基本域，取正向边界：底部圆弧从 \(\rho\) 经 \(\mathrm{i}\) 走到 \(\rho+1\)，顶部由右向左。周期性使竖边的 \(f'/f\) 积分相消；顶部积分趋于 \(-2\pi\mathrm{i}v_\infty(f)\)。
把底弧分为左半弧 \(L:\rho\to\mathrm{i}\) 与右半弧 \(R:\mathrm{i}\to\rho+1\)。映射 \(S\) 将 \(L\) 映成反向的 \(R\)。微分变换律给
\(\mathrm{d}\log f(S\tau)=\mathrm{d}\log f(\tau)+k\,\mathrm{d}\tau/\tau\)。故两弧积分之和等于 \(-k\int_L\mathrm{d}\tau/\tau\)。左弧的辐角从 \(2\pi/3\) 减至 \(\pi/2\)，结果为 \(k\pi\mathrm{i}/6\)。辐角原理除以 \(2\pi\mathrm{i}\) 得内部零极点总数等于 \(k/12-v_\infty\)。
若边界有零极点，先删去其小圆邻域，并将成对边界上的删去部分配成映射下的对应段。\(f'/f=v_p/(\tau-p)+O(1)\)，所以小弧积分的极限等于其角度占 \(2\pi\) 的比例乘 \(2\pi\mathrm{i}v_p\)。普通边界点的两份半圆合成整圆；\(\mathrm{i}\) 只占角度 \(\pi\)；\(\rho,\rho+1\) 两角合计 \(2\pi/3\)。于是分别得到系数一、二分之一、三分之一。被删除的小段对权修正积分的贡献趋零，顶部仍用同一 \(q\) 展开估计。取小圆半径趋零即得所示公式。
\end{proof}
这一结果称为\term[jia-gongshi]{价公式}[valence formula]。Zagier 的《Elliptic Modular Forms and Their Applications》\cite[Proposition 2, pp. 9--10]{Zagier2008Elliptic}可用于核对上述固定点权重；本书在围道计算中显式保留了边界方向。

\subsection{维数论证}
\label{b3:subsec:230503}

\begin{theorem}\label{b3:thm:23-modular-ring}
全模群的分次模形式环为 \(M_* =\mathbb C[E_4,E_6]\)，两个生成元代数独立。对偶数 \(k\ge0\)，
\[
\dim M_k=\begin{cases}\lfloor k/12\rfloor,&k\equiv2\pmod{12},\\
\lfloor k/12\rfloor+1,&k\not\equiv2\pmod{12}.
\end{cases}
\]
负权及奇权空间均为零。
\end{theorem}
\begin{proof}
价公式各项非负，故负权无非零形式；奇权已由 \(-I\) 排除。权零形式无零极点，或用紧化基本域上的最大值原理，可知为常数。权二时 \(v_i\ge1\)，其贡献至少 \(1/2>2/12\)，故空间为零。
因 \(\Delta\) 内部无零点且尖点一阶消失，除以 \(\Delta\) 给出 \(S_k=\Delta M_{k-12}\)。每个偶数 \(k\ge4\) 都可写成 \(4a+6b\)，其中 \(a,b\ge0\)：\(k\equiv0\pmod4\) 时取 \(b=0\)，\(k\equiv2\pmod4\) 时取 \(b=1\)。从 \(f\in M_k\) 减去其常数项乘 \(E_4^aE_6^b\)，余项属于 \(S_k\)。对权归纳，得到所有形式均为两生成元的多项式。每次常数项映射满射到 \(\mathbb C\)，故 \(\dim M_k=1+\dim M_{k-12}\)；结合 \(0,2,4,6,8,10\) 六个初值即得公式。
为证代数独立，先把任何多项式关系按权分解。对固定 \(\tau\)，代入无限多个矩阵 \(\left(\begin{smallmatrix}1&0\\n&1\end{smallmatrix}\right)\)，得到不同权部分乘 \((n\tau+1)^k\) 的多项式恒等式；不同次数的系数必须各为零。因此只需排除同权关系。同权单项式约去一个共同因子后，成为 \(E_4^3/E_6^2\) 的多项式关系。这个亚纯函数非常数，因为两者常数项均一而一次系数分别为 \(720,-1008\)。非常值亚纯函数在没有极点的某小圆盘上像含开集，不可能满足非零常系数多项式，故关系为零。
\end{proof}
例如 \(M_8\) 一维，常数项为一的 \(E_8\) 必等于 \(E_4^2\)。比较 Fourier 系数即得到约数和卷积恒等式。有限维数把一个无穷级数恒等式缩减为有限个系数的核对；减去已匹配项后，价公式保证余项必须为零。

\ProseParagraphBreak

模形式环还允许构造一个保持模性的微分运算。普通导数的变换律含有额外项，必须先计算这项，而不能直接把导数称为高二权形式。记 \(D=(2\pi\mathrm{i})^{-1}\mathrm{d}/\mathrm{d}\tau=q\mathrm{d}/\mathrm{d}q\)，并令
\[
 P(\tau)=\frac{D\Delta(\tau)}{\Delta(\tau)}.
\]
判别式在上半平面无零点，故 \(P\) 全纯。尖点处 \(\Delta=q(1-24q+\cdots)\)，所以 \(P=1-24q+O(q^2)\)。对 \(\Delta(\gamma\tau)=(c\tau+d)^{12}\Delta(\tau)\) 求对数导数，链式法则给
\[
 P(\gamma\tau)=(c\tau+d)^2P(\tau)
       +\frac{12c(c\tau+d)}{2\pi\mathrm{i}}.
\]
额外的加项说明它不是权二模形式；这与 \(M_2=0\) 完全一致。它的任务是抵消其他函数导数中的同类加项。
\ProseParagraphBreak
对权 \(k\) 形式 \(f\) 求导，得到
\[
 (Df)(\gamma\tau)=(c\tau+d)^{k+2}Df(\tau)
       +\frac{kc}{2\pi\mathrm{i}}(c\tau+d)^{k+1}f(\tau).
\]
故 \(D_kf=Df-(k/12)Pf\) 满足权 \(k+2\) 的变换律。其 \(q\) 展开没有负幂；若 \(f\) 是尖点形式，两个乘积也仍无常数项。因此 \(D_k\) 把 \(M_k\) 映入 \(M_{k+2}\)，并保持尖点子空间。这是对普通导数的一项具体修正，全部常数都由权十二的判别式决定。
\ProseParagraphBreak
现在用已经证明的低权维数即可得微分恒等式。\(D_4E_4\in M_6\)，其常数项为 \(-1/3\)，所以 \(D_4E_4=-E_6/3\)。同理 \(D_6E_6\in M_8=\mathbb C E_4^2\)，常数项为 \(-1/2\)。于是
\[
 DE_4=\frac{PE_4-E_6}{3},\qquad
 DE_6=\frac{PE_6-E_4^2}{2}.
\]
再对 \(P\) 的变换律求导，可得 \(DP-P^2/12\) 为权四模形式。具体而言，令 \(J=c\tau+d\)，导数变换中的额外两项是 \(2cJ^3P/(2\pi\mathrm{i})\) 与 \(12c^2J^2/(2\pi\mathrm{i})^2\)；展开 \(P(\gamma\tau)^2/12\) 恰给相同两项，因而完全相消。其常数项为 \(-1/12\)，故
\[
 DP=\frac{P^2-E_4}{12}.
\]
这些式子展示有限维性怎样将微分信息转为代数关系。它们不依赖尚未证明的无限乘积公式，也没有从前几项相同就直接宣称两函数相同；关键是先证明差属于一个已知的一维模形式空间。
\ProseParagraphBreak
也可反过来从这组微分方程递推系数。若写 \(P=1+\sum p_nq^n\)、\(E_4=1+\sum a_nq^n\)、\(E_6=1+\sum b_nq^n\)，每个 \(D\) 把第 \(n\) 项乘以 \(n\)，乘积的第 \(n\) 项则是有限卷积。给定前面的系数后，方程约束后续系数，而 Eisenstein 级数的显式约数和满足这些约束。这是可计算的一致性检验，但如果要据此从零定义三个函数，仍须另外证明所得形式级数的收敛；本章已有全纯构造，故无需倒置这一步。
